\chapter{Elementary Set Theory as Applied Logic}
\label{ch:set-theory}

This closing chapter treats membership, subset relations,
comprehension over typed domains, unions, intersections, and
complements as applied instances of the logic already developed,
rather than as a separate subject. Russell's paradox is used to draw
a sharp distinction between an inconsistent specification and a
merely unproved proposition.

\section{Membership and Comprehension}

In Lean, \texttt{Set alpha} is not a new kind of object requiring
new logic — it is literally \texttt{alpha → Prop}, a predicate on
\texttt{alpha} wearing set-theoretic notation. Comprehension,
\texttt{\{x | p x\}}, builds a set directly from a predicate
\texttt{p}, and membership, $x \in s$, unfolds to nothing more
than applying that predicate: \texttt{s x}. Every subset proof is
therefore already a proof about an implication, provable with
exactly the tactics Chapter~\ref{ch:natural-deduction} introduced.

\begin{experiment}{A subset proof via \texttt{intro}/\texttt{exact}}
\begin{lean}
example : {x : Nat | x < 3} ⊆ {x : Nat | x < 5} := by
  intro x hx
  simp only [Set.mem_setOf_eq] at hx ⊢
  omega
\end{lean}
Unfolding $\subseteq$ at the definitional level, the goal is exactly
an implication — $x < 3 \to x < 5$ — and \texttt{intro x hx}
discharges it the same way \texttt{intro} discharged any other
implication in Chapter~\ref{ch:natural-deduction}. The set-theoretic
vocabulary changed; the proof method did not.
\end{experiment}

\section{Unions, Intersections, Complements}

De Morgan's laws hold for sets in exactly the shape they held for
propositions in Chapter~\ref{ch:propositions}, because membership in
a union, an intersection, or a complement is defined to unfold to
$\lor$, $\land$, and $\lnot$ respectively — proving the set version
is proving the propositional version, one membership check away.

\begin{experiment}{De Morgan for sets, alongside De Morgan for propositions}
\begin{lean}
example {alpha : Type} (A B : Set alpha) :
    (A ∪ B)ᶜ = Aᶜ ∩ Bᶜ := by
  ext x
  simp only [Set.mem_compl_iff, Set.mem_union, Set.mem_inter_iff]
  exact Iff.intro
    (fun h => And.intro (fun hA => h (Or.inl hA)) (fun hB => h (Or.inr hB)))
    (fun h => fun hab => hab.elim h.1 h.2)
\end{lean}
Set this beside \texttt{demorgan\_or} from
Chapter~\ref{ch:propositions}: after \texttt{ext x} reduces the
set equation to a statement about the single element \texttt{x}, and
the three \texttt{simp} lemmas unfold membership, the two proof terms
are the identical shape — \texttt{And.intro} built from two
functions composing with \texttt{Or.inl}/\texttt{Or.inr}, and
\texttt{.elim} case-splitting an \texttt{Or} in the other direction.
Nothing new was proved; the same fact was proved again in
different notation, and the proof noticed.
\end{experiment}

\section{Russell's Paradox: An Inconsistency, Not an Open Question}

\begin{experiment}{Why unrestricted comprehension fails}
\begin{lean}
-- This does NOT typecheck — Lean rejects it before there is any
-- proposition left to ask "true or false" about:
def russellSet {alpha : Type} : Set alpha := {s | s ∉ s}
-- Elaboration fails: `s ∉ s` needs `s : alpha` (to be a candidate
-- element of the set being built) and also `s : Set alpha` (to be
-- one of the sets `s` itself ranges over) at once — that is, it
-- needs `alpha = Set alpha`. No fixed Lean `Type` can equal the
-- type of its own subsets; the universe hierarchy is stratified
-- precisely to rule this out.
\end{lean}
Naive set theory's version of this — ``let $R$ be the set of all
sets that do not contain themselves; is $R \in R$?'' — treats the
attempted definition as though it succeeds, and only then discovers
a contradiction from asking whether $R \in R$. Lean's type theory
never reaches that question: the definition of \texttt{russellSet}
fails to elaborate at all, because $s \notin s$ is not
well-typed for an arbitrary \texttt{s : alpha} in the first place.
The paradox is blocked structurally, one layer before it could ever
become a true-or-false question.
\end{experiment}

\section{Closing the Introductory Sequence}

This is the last chapter of the introductory sequence.
\textit{Inference Under Constraint} continues from here into
non-classical logic, modality, and a broader admissibility-theoretic
program; nothing in this book depends on that continuation, and a
reader who stops here has a complete, self-contained introduction to
logic through Lean.

\section{Exercises}

\begin{exercise}
Prove $A \cap (B \cup C) = (A \cap B) \cup (A \cap C)$ for sets over a
common type.
\end{exercise}
